\chapter{Considerações Finais}
\label{chp:consideracoes_finais}

O Projeto Pedagógico do Curso de \ppccurso da Universidade Federal de Uberlândia reflete o compromisso da instituição com a formação de profissionais éticos, tecnicamente qualificados, aplicados e socialmente responsáveis. Estruturado em consonância com o Catálogo Nacional de Cursos Superiores de Tecnologia (CNCST, 4ª edição)~\cite{mec_cncst_2024}, com as Diretrizes Curriculares Nacionais Gerais para a Educação Profissional e Tecnológica~\cite{cne_cp_1_2021} e com o Plano Institucional de Desenvolvimento e Expansão (PIDE) da UFU~\cite{ufu_consun_31_2022}, este PPC apresenta uma proposta formativa fundamentada em competências, interdisciplinaridade, integração teoria-prática e valorização da diversidade humana e cultural.

A matriz curricular, os objetivos do curso, o perfil do egresso, as estratégias de ensino e os mecanismos de avaliação foram definidos para garantir uma formação alinhada aos desafios da automação industrial contemporânea — instrumentação, controle, redes industriais e, na ênfase própria deste curso, robótica e inteligência artificial aplicada (\ppcenfasecurricular) — e às demandas do setor produtivo. O curso busca oferecer ao(à) estudante um percurso formativo aplicado, crítico, flexível e capaz de dialogar com a complexidade da indústria automatizada, da ciência e da vida em sociedade.

Este PPC também apresenta a proposta de ABI com a Engenharia de Automação, Robótica e Inteligência Artificial (\autoref{chp:abi}), apoiada na identidade dos cinco primeiros períodos das matrizes. O ingresso, a escolha da trajetória e qualquer continuidade em uma segunda graduação ainda dependem de ato institucional específico; este PPC não cria transferência, permanência de vínculo ou dupla diplomação automática.

Este documento expressa também a valorização do diálogo entre discentes, docentes, técnicos(as) e gestores(as), bem como o papel estratégico da universidade pública no desenvolvimento regional e nacional. O \ppccurso da FEELT/UFU reafirma, assim, sua missão de formar tecnólogas e tecnólogos comprometidos com a excelência técnica, a segurança e a sustentabilidade dos processos industriais, a justiça social e a construção de um futuro sustentável.
